\documentclass[a4paper,11pt]{article}
\usepackage{exam-style-341}
\examinehead{341 农业综合知识三 · 程序设计分章练习}{第06章参考答案}
\begin{document}

\unitanswerscover{341 农业综合知识三 · 程序设计分章练习}{第06章参考答案}{数组：逐题答案、解析与综合题完整参考实现；判断题均给出改正后的表述，程序阅读与参考程序均已用 \code{gcc -std=c99} 实跑验证。}

\answerpart{一、单项选择题（10 题，每题 2 分，共 20 分）}

\q C。\begin{analysis}\code{int a[5]=\{1,2\};} 只给前两个元素赋初值，其余元素自动初始化为 0，因此 \code{a[3]} 的值为 0。A、B 分别是 \code{a[0]}、\code{a[1]} 的值；D 错，部分初始化时未赋值元素有确定的初值 0，不是随机值。\end{analysis}

\q A。\begin{analysis}数组元素下标从 0 开始，长度为 10 的数组 \code{a[10]} 的合法下标是 0\textasciitilde 9；\code{a[10]} 已越界。B、D 从 1 开始计数，C 把长度 10 当成了合法下标。\end{analysis}

\q B。\begin{analysis}\code{int a[]=\{1,2,3,4\};} 省略方括号中的长度时，由初值个数自动确定数组长度为 4。C 错在把结束符规则套用到整型数组；A、D 与定义不符。\end{analysis}

\q B。\begin{analysis}二维数组按行优先存放：\code{a[0][0]} 是第 1 个，\code{a[0][3]} 是第 4 个，\code{a[1][0]} 是第 5 个，\code{a[1][1]} 是第 6 个，\code{a[1][2]} 是第 7 个。也可用公式 $i\times4+j+1=1\times4+2+1=7$ 计算。\end{analysis}

\q B。\begin{analysis}用字符串常量初始化字符数组时，末尾自动加上字符串结束符 \code{'\textbackslash 0'}，所以 \code{char s[]="abc";} 的长度为 $3+1=4$，\code{sizeof(s)} 为 4。A 只数了 3 个字符；C、D 与规则不符。\end{analysis}

\q A。\begin{analysis}\code{char s[10]="abc";} 数组长度为 10，但字符串有效字符只有 3 个，\code{strlen} 数到第一个 \code{'\textbackslash 0'} 为止，结果为 3；\code{sizeof(s)} 才是 10。B 错在把 \code{'\textbackslash 0'} 计入 strlen；C 是 sizeof 的结果；D 错，strlen 的结果是确定的。\end{analysis}

\q B。\begin{analysis}字符数组不能用 \code{=} 整体赋值（A、C、D 都试图对数组名赋值，其中 C 还缺少定义形式，语法错误），正确做法是定义时初始化或用 \code{strcpy} 复制。D 中 \code{s="hello"+1} 既是对数组名赋值，又把字符串指针当作算术运算对象，均错误。\end{analysis}

\q B。\begin{analysis}数组名作函数实参时，传递的是数组首元素的地址（即数组名所代表的地址），形参 \code{int a[]} 在函数中按指针处理。因此函数内对 \code{a[i]} 的修改会作用于实参数组。A、D 错在认为传递整个数组；C 错在数组长度不会自动传递。\end{analysis}

\q A。\begin{analysis}循环 \code{i} 取 1、2、3，累加 \code{a[1]+a[2]+a[3]=4+3+2=9}；\code{i=0} 和 \code{i=4} 不参与。B 误把 \code{a[4]} 也算入（9+1=10 或误算为 12）；C 是全部元素之和 15；D 只是前三个元素之和。\end{analysis}

\q C。\begin{analysis}\code{strcmp("abc","abd")} 按字符逐个比较，前两个字符相同，第三个字符 \code{'c'} 的 ASCII 码（99）小于 \code{'d'}（100），所以返回负数（具体值由实现决定，不一定是 $-1$）。A 表示第一个字符串大，B 表示相等，D 把返回值误当成固定值。\end{analysis}

\answerpart{二、多项选择题（5 题，每题 3 分，共 15 分）}

\q ABC。\begin{analysis}A 完全初始化正确；B 部分初始化，后两个元素自动为 0，正确；C 省略长度由初值个数确定长度 3，正确；D 给长度为 3 的数组提供 4 个初值，超出数组范围，属于越界初始化，编译会报错或产生未定义行为。\end{analysis}

\q ABC。\begin{analysis}A 正确：用字符串常量初始化字符数组时自动加 \code{'\textbackslash 0'}；B 正确：部分初始化时未赋值元素为 0，即 \code{'\textbackslash 0'}；C 正确：strlen 不把 \code{'\textbackslash 0'} 计入长度；D 错：字符数组不能直接用 \code{=} 赋值字符串（只能在定义时用字符串常量初始化，之后要用 strcpy）。\end{analysis}

\q ABC。\begin{analysis}A、B 正确：形参可以写成 \code{int a[]}，也可以写成 \code{int *a}，二者等价；C 正确：实参传递的是数组首地址；D 错：形参 a 在函数内是指针，\code{sizeof(a)} 得到指针的大小（与平台有关），不是整个数组的字节数，所以求数组长度必须另设长度参数。\end{analysis}

\q ABC。\begin{analysis}strcpy（复制）、strcmp（比较）、strcat（连接）都是 \code{<string.h>} 中的字符串处理函数；printf 是格式化输出函数，属于 \code{<stdio.h>}，不是字符串处理函数。\end{analysis}

\q ABC。\begin{analysis}A 正确：2 行 3 列共 6 个元素；B 正确：可省略第一维、由初值确定，第二维不能省；C 正确：二维数组按行优先顺序连续存放；D 错：第二维长度不能省略，否则编译器无法确定行的大小。\end{analysis}

\answerpart{三、判断题（5 题，每题 2 分，共 10 分）}

\q √。\begin{analysis}数组下标从 0 开始，长度为 10 的数组合法下标为 0\textasciitilde 9，\code{a[10]} 是越界访问。\end{analysis}

\q ×。改正：\code{char s[]="abc";} 的长度是 4，末尾含 \code{'\textbackslash 0'}；\code{char s[]=\{'a','b','c'\};} 的长度是 3，不含 \code{'\textbackslash 0'}，不能直接按字符串使用。\begin{analysis}用字符串常量初始化时，编译器自动在末尾加结束符；逐个字符初始化时不会自动加，除非初值个数少于数组长度（未赋值元素为 0）。\end{analysis}

\q √。\begin{analysis}数组名代表数组首元素的地址，是地址常量，不能赋值、不能自增；要移动访问应使用指针变量，如 \code{int *p=a; p++;}。\end{analysis}

\q ×。改正：判断两个字符串是否相等应使用 \code{strcmp(s1,s2)==0}，不能直接用 \code{==} 比较字符数组名。\begin{analysis}数组名在比较时会被转换为地址，\code{s1 == s2} 比较的是两个地址是否相同，而不是字符串内容是否相同。\end{analysis}

\q ×。改正：字符数组不能用 \code{=} 整体赋值，只能在定义时用字符串常量初始化，或在定义后用 \code{strcpy} 复制。\begin{analysis}\code{char s[10]; s="hello";} 中 s 是数组名（地址常量），不能被赋值，编译会报错。正确写法是 \code{char s[10]="hello";} 或 \code{strcpy(s,"hello");}。\end{analysis}

\answerpart{四、填空题（10 空，每空 1 分，共 10 分）}

\q （1）0；（2）0。\begin{analysis}部分初始化时，未显式赋值的元素自动初始化为 0，因此 \code{a[2]}、\code{a[4]} 均为 0。\end{analysis}

\q （3）12；（4）7。\begin{analysis}$3\times4=12$ 个元素；按行优先，\code{a[1][2]} 前面有第 0 行的 4 个元素和第 1 行的 2 个元素共 6 个，所以它是第 7 个元素。\end{analysis}

\q （5）5；（6）6。\begin{analysis}"China" 有 5 个字符，strlen 不计结束符，结果为 5；数组长度要容纳 5 个字符加 1 个 \code{'\textbackslash 0'}，sizeof 为 6。\end{analysis}

\q （7）\code{strcpy}；（8）\code{strcat}。\begin{analysis}字符串复制用 strcpy，字符串连接用 strcat，二者都要求目标数组有足够空间。\end{analysis}

\q （9）首地址。\begin{analysis}数组名作函数实参时传递数组首元素的地址，形参按指针接收，函数内可修改实参数组的元素。\end{analysis}

\q （10）5。\begin{analysis}省略长度时由初值个数确定数组长度，共 5 个初值，长度为 5。\end{analysis}

\answerpart{五、程序阅读题（2 题，每题 5 分，共 10 分）}

\q 运行结果为：
\begin{lstlisting}
6
\end{lstlisting}
\begin{analysis}\code{int a[8]=\{1,2,3\};} 只给前 3 个元素赋值，其余 5 个元素自动为 0。循环把 8 个元素全部累加：$1+2+3+0\times5=6$，所以输出 6。本题的陷阱是误以为未赋值元素是随机值，或误以为数组长度是 3。\end{analysis}

\q 运行结果为：
\begin{lstlisting}
ChinaBeijing 12 20
\end{lstlisting}
\begin{analysis}初始 \code{s1} 为 "China"（占 20 字节），\code{s2} 为 "Beijing"；\code{strcat(s1,s2)} 把 "Beijing" 连接到 "China" 末尾，得到 "ChinaBeijing"（12 个字符）；\code{strlen(s1)} 返回字符个数 12，不计 \code{'\textbackslash 0'}；\code{sizeof(s1)} 返回数组 s1 的字节数 20，与内容无关。\end{analysis}

\answerpart{六、简答题（3 题，每题 5 分，共 15 分）}

\q 参考答案要点：
（1）定义形式：\code{类型说明符 数组名[常量表达式];}，如 \code{int a[5];}，方括号内表示元素个数，下标从 0 开始，合法下标为 0\textasciitilde 4（2 分）。
（2）初始化规则：可全部初始化 \code{int a[5]=\{1,2,3,4,5\};}；可部分初始化 \code{int a[5]=\{1,2\};}，未赋值元素自动为 0；可省略长度 \code{int a[]=\{1,2,3\};}，由初值个数确定长度为 3；初值个数不能超过数组长度（2 分）。
（3）部分初始化：只给数组前若干个元素赋初值，其余元素自动初始化为 0。越界初始化：给出的初值个数超过数组长度，如 \code{int a[3]=\{1,2,3,4\};}，属于语法错误或未定义行为（1 分）。

\q 参考答案要点：
（1）字符数组是元素类型为 char 的数组，可以只存放若干字符而不含结束符，如 \code{char c[5]=\{'C','h','i','n','a'\};}；字符串是以 \code{'\textbackslash 0'} 结尾的字符序列，用字符串常量初始化时末尾自动加 \code{'\textbackslash 0'}，如 \code{char s[]="China";}（2 分）。
（2）strlen(s)：从 s 开始数到第一个 \code{'\textbackslash 0'} 的字符个数，不计 \code{'\textbackslash 0'}；对 \code{char s[]="China";} 返回 5（1.5 分）。
（3）sizeof(s)：s 所占的存储字节数；若 s 是数组则包含 \code{'\textbackslash 0'}，对 \code{char s[]="China";} 返回 6；若 s 是函数形参（指针）则返回指针大小，与字符串内容无关（1.5 分）。

\q 参考答案要点：
（1）形参写法：可以写成 \code{int a[]}，也可以写成 \code{int *a}，二者等价，a 在函数内都按指针处理（2 分）。
（2）传递内容：实参写数组名（或数组首元素地址），传递的是数组首元素的地址，不是整个数组的副本；函数内通过 \code{a[i]} 或 \code{*(a+i)} 可以访问并修改实参数组元素（2 分）。
（3）原因：形参 a 实质上是指针变量，\code{sizeof(a)} 求的是指针变量本身的大小（如 8 字节或 4 字节），不是实参数组的字节数，因此函数内无法用 sizeof 求数组长度，必须把数组长度作为另一个参数传入（1 分）。

\answerpart{七、综合题（1 题，共 20 分）}

\q 参考程序如下：
\begin{lstlisting}[language={[custom]C}]
#include <stdio.h>
int main(void)
{
    int a[3][4], i, j, max, r = 0, c = 0;
    for (i = 0; i < 3; i++)
        for (j = 0; j < 4; j++)
            scanf("%d", &a[i][j]);
    max = a[0][0];
    for (i = 0; i < 3; i++)
    {
        int rowsum = 0;
        for (j = 0; j < 4; j++)
        {
            rowsum += a[i][j];
            if (a[i][j] > max) { max = a[i][j]; r = i; c = j; }
        }
        printf("row%d=%d\n", i + 1, rowsum);
    }
    printf("max=%d,row=%d,col=%d\n", max, r + 1, c + 1);
    for (j = 0; j < 4; j++)
    {
        for (i = 0; i < 3; i++)
            printf("%d ", a[i][j]);
        printf("\n");
    }
    return 0;
}
\end{lstlisting}
\begin{analysis}用输入 \code{1 2 3 4 5 6 7 8 9 10 11 12} 实跑，输出为：
\begin{lstlisting}
row1=10
row2=26
row3=42
max=12,row=3,col=4
1 5 9
2 6 10
3 7 11
4 8 12
\end{lstlisting}
三行之和分别为 $1+2+3+4=10$、$5+6+7+8=26$、$9+10+11+12=42$；最大值为 12，位于第 3 行第 4 列；转置输出时外层循环控制列 j、内层循环控制行 i，得到 $4\times3$ 矩阵。\end{analysis}
\begin{analysis}评分要点：正确定义 $3\times4$ 二维数组（3 分）；用二重循环和 \code{scanf} 读入全部 12 个元素，\code{\&a[i][j]} 地址形式正确（4 分）；正确计算并输出每行之和（3 分）；正确找出最大值并记录行、列下标（4 分）；输出时行号、列号从 1 开始（2 分）；转置输出时行列循环次序正确（4 分）。\end{analysis}
\begin{analysis}易错提示：\code{scanf} 中漏写 \code{\&}；把行列循环写反导致按行输出而非转置输出；记录最大值时忘记同步更新行、列下标；输出行号时直接输出下标 i 而没有加 1；把行和变量 \code{rowsum} 定义在行循环之外导致累加不归零。\end{analysis}

\end{document}
